\chapter{Edema (Fluid Retention)}

\section*{Definition}
Edema is the accumulation of excess fluid in the interstitial space of tissues, causing visible swelling. It can be localized (venous insufficiency, lymphatic obstruction) or generalized (heart failure, kidney disease, liver disease).

\section*{What Happens in Your Body}
Edema arises from an imbalance in Starling forces: (1) increased capillary hydrostatic pressure (heart failure, venous insufficiency); (2) decreased plasma oncotic pressure (hypoalbuminemia); (3) increased capillary permeability (inflammation, infection, allergic reaction); (4) lymphatic obstruction (lymphedema).

\section*{Causes}
\begin{itemize}
\item Chronic venous insufficiency (most common cause of bilateral leg edema)
\item Congestive heart failure
\item Cirrhosis with hypoalbuminemia
\item Nephrotic syndrome
\item Acute or chronic kidney disease
\item Deep vein thrombosis (unilateral edema)
\item Lymphedema (cancer treatment, infection, congenital)
\item Medications (calcium channel blockers, NSAIDs, corticosteroids)
\item Pregnancy (physiologic edema)
\item Idiopathic edema
\end{itemize}

\section*{When to See a Doctor}
\begin{itemize}
\item Visible swelling of one or both legs, ankles, feet
\item Pitting edema on examination
\item Rapid weight gain over days
\item Swelling accompanied by shortness of breath (possible heart failure)
\item Unilateral leg swelling with pain (possible DVT)
\item Periorbital edema (suggesting kidney disease)
\item Ascites or abdominal distension
\item Skin changes: thinning, dry, ulcerated skin
\end{itemize}

\section*{Related Symptoms}
\begin{itemize}
\item Shortness of breath (if pulmonary edema)
\item Ascites (abdominal fluid)
\item Weight gain (fluid retention)
\item Orthopnea (difficulty breathing when lying flat)
\item Paroxysmal nocturnal dyspnea
\item Skin ulcers (from chronic venous stasis)
\item Cellulitis (infection of edematous skin)
\end{itemize}

\section*{Self-Care / Home Management}
\begin{itemize}
\item Elevate the legs above heart level several times daily
\item Reduce dietary sodium to less than 2 g per day
\item Wear compression stockings if venous insufficiency
\item Avoid prolonged sitting or standing
\item Stay hydrated
\item Limit alcohol consumption
\item Monitor daily weight and ankle circumference
\end{itemize}

\section*{How Your Doctor Will Evaluate You}
\begin{itemize}
\item History and physical examination
\item Basic metabolic panel and serum albumin
\item Urinalysis for proteinuria
\item Brain natriuretic peptide (BNP) if heart failure suspected
\item Liver function tests if cirrhosis suspected
\item Venous duplex ultrasound if DVT suspected
\item Chest X-ray if pulmonary edema suspected
\end{itemize}

\section*{Treatment Options}
\begin{itemize}
\item Treat underlying cause (heart failure, liver disease, kidney disease)
\item Diuretic therapy (loop diuretics: furosemide)
\item Sodium restriction
\item Compression therapy for venous insufficiency or lymphedema
\item ACE inhibitors for nephrotic syndrome
\item Discontinue contributing medications
\item Leg elevation and mobilization
\end{itemize}

\section*{References}
\begin{thebibliography}{99}
\bibitem{cho2016}
Cho S, Atwood JE. Peripheral edema. \textit{Am J Med}. 2002;113(7):580--586.
\bibitem{trayes2013}
Trayes KP, Studdiford JS, Pickle S, Tully AS. Edema: diagnosis and management. \textit{Am Fam Physician}. 2013;88(2):102--110.
\bibitem{braunwald2008}
Braunwald E. Heart failure. In: \textit{Harrison's Principles of Internal Medicine}. 17th ed. McGraw-Hill; 2008.
\bibitem{ellison2000}
Ellison DH. Edema. In: \textit{Brenner \& Rector's The Kidney}. 6th ed. WB Saunders; 2000.
\bibitem{yosipovitch2004}
Yosipovitch G, Sackett-Lundeen L, Goon A, et al. Circadian and ultradian rhythms of body weight and peripheral edema. \textit{Chronobiol Int}. 2004;21(3):455--465.
\bibitem{gorski2015}
Gorski LA, Hadaway L, Hagle ME, et al. Infusion therapy standards of practice. \textit{J Infus Nurs}. 2016;39(1 suppl):S1--S159.
\end{thebibliography}



\clearpage
